%! Author = sriadityavedantam
%! Date = 3/29/26

The content in this chapter is things to know by heart.
We will not be going back and explaining the content discussed in this chapter.

\section{Logic}

For those coming from a pure symbolic proofs-based class, this text will definitely be a bit striking, as I do not like using symbols every time they can be used.
It is easier to convey thoughts by just using words and to depict very slight meanings that may not be robotic.
It is definitely not impossible to do the mental conversion into symbolic language.
However, the way I learned proofs was to use more words than symbols.
As a matter of fact, some classes may even deduct points for the overuse of symbols, and I have heard this tale through and through from many people.
So take what you will, but I hope this will create some change.
If there is one thing to take away from this section, it is that there is nothing ever wrong with using words over symbols, while there is the vice versa.

Let $P$ and $Q$ be statements.
It should have been discussed in a proof class the difference between statements, questions, and commands.

\begin{definition}[Basic Logical Statements]
    \textbf{``$P$ and $Q$''}
    This is true if and only if $P$ and $Q$ are both true.
    This is denoted by $\land$.

    \textbf{``$P$ or $Q$''}
    This is true in all cases where at least one of $P$ or $Q$ is true, and false only when they are both false.
    This is denoted by $\lor$.

    \textbf{``$P$ implies $Q$''}
    We use implications to show that if $P$ holds, then $Q$ follows.
    For example, we usually write this in English as ``If $P$, then $Q$.''
    This means that if $P$ is true, then $Q$ will also happen.
    This is true in $3/4$ of the possible outcomes.
    Namely, it is true when both $P$ and $Q$ are true, when both are false, and when $P$ is false but $Q$ is true.
    It is false only when $P$ is true and $Q$ is false.
    A false premise always makes the implication true.
    Implications are denoted by $\implies$.

    \textbf{``$P$ if and only if $Q$''}
    This is called a biconditional, or an equivalence statement.
    This is short for saying ``$P$ implies $Q$ and $Q$ implies $P$.''
    This is denoted by $\iff$.

    \textbf{``It is not the case that $P$''}
    This is true if and only if $P$ is false.
    This is also called negation.
\end{definition}